\subsection{Policy routes, pinned memories, and bounded delivery}
\label{sec:policy_delivery_update}

The September 23 source update makes explicit routing part of candidate generation. A recall-policy entry has one of three modes: \texttt{normal}, \texttt{parent\_only}, or \texttt{disabled}. Parent-only routes require their approved direct lexical or intent condition, while disabled entries are removed from ordinary deliberate retrieval. The new behavior lets a normal entry carry multi-term trigger groups and explicit token aliases. When every term of any group occurs in the alias-enriched query and no configured blocking phrase matches, the entry receives a policy seed of strength $3$ and a lexical seed of at least $3$. Without a trigger, it retains ordinary retrieval eligibility. This mechanism is a hand-authored routing rule; matching a reviewed spelling or token group does not establish semantic generalization \citep{brainrecall20260923,recallpolicy20260923}.

Policy precedence also required a repair. The original promotion placed a policy seed above the initial BM25 peak by an additive margin. Later multiplicative name-coverage and question-coverage boosts could move an ordinary result above it. The revised default BM25-first path lifts policy candidates again after those boosts. Writing $b_j$ for the adjusted non-policy scores and $s_i$ for a policy strength, the implemented lift is
\[
 b_i^{\mathrm{policy}}
 = \max_{j\notin P} b_j + 2 + \frac{s_i}{s_{\max}},
\]
with a fallback to the initial BM25 peak when there is no non-policy score. This restores precedence at that stage while retaining policy-strength order. It is a scoped implementation property: \texttt{BRAIN\_RECALL\_BLEND=1} selects the legacy blend, and the later optional status-order pass can rerank candidates; that pass defaults to off. A change to a subsequent ranking stage therefore requires another precedence check \citep{brainrecall20260923}.

Pins have an opportunity cost. Deliberate recall defaults to eight displayed rows, and a pinned entry occupies an ordinary row in that allowance. Moreover, the seed is inserted before graph propagation and seed-based semantic expansion. The implementation consequently does not guarantee that removing the pinned row would recover an otherwise identical result list. Such preservation must be measured on a stated query set. The pin switch, \texttt{BRAIN\_RECALL\_PINS=0}, disables the new normal-entry promotions. Independently, displayed descriptions now prefer the memory's own frontmatter, read from a bounded prefix, and default to a 320-character excerpt with an ellipsis when shortened \citep{brainrecall20260923}.

Delivery imposes further limits beyond retrieval rank. The dispatcher retains its 9,800-character allocation, whereas owner-history delivery defaults to a 9,000-UTF-8-byte budget after accounting for existing context and its history notice. These quantities are not model-token counts. The new dispatcher can skip ten designated operational hooks for recognized headless workers when its marker is present, and avoids launching a hook whose explicit enabling file is absent. These controls specify conditional source behavior; source presence does not establish which marker or settings a deployed session used \citep{dispatcher20260923,ownerdelivery20260923}.

The evidence boundary remains explicit. Owner-history state records a context hash and byte length only after successful stdout writing and flushing, labeling the receipt \texttt{emitted\_stdout\_not\_receiver\_verified}. Retrieval and emission therefore require separate evidence from receipt, action, and outcome. Similarly, the revised regression gate accepts selected growing-input receipt differences within a default three-hour window, while still rejecting other source-field drift. A passing freshness check under that rule is compatible with limited input growth; it is not an assertion that every source field is identical to the build snapshot \citep{ownerdelivery20260923,braingate20260923}.
